% =====================================================================
% Frente 5 --- Influenciadores (Ravi)
% =====================================================================
\section{Frente 5 --- Influenciadores}

\noindent\textbf{Responsável:} Ravi.
\textbf{No relatório:} problema aplicado e resultados dos influenciadores; referências em BibTeX.
\textbf{Na Etapa 1:} descrição formal do problema aplicado.

\subsection{O que foi implementado}

{\small\begin{tabularx}{\linewidth}{@{}L{3.9cm} L{4.7cm} Y L{1.9cm}@{}}
\toprule
Função & Arquivo & Complexidade & Situação \\
\midrule
\codigo{limiar\_epidemico} & \codigo{aplicacao/influencia.py} & $O(n)$~\cite{newman2002} & \feito \\
\codigo{simular\_sir} (uma rodada) & \codigo{aplicacao/influencia.py} & $O(n + m)$ no pior caso, com $\gamma = 1$ & \feito \\
\codigo{poder\_de\_espalhamento} & \codigo{aplicacao/influencia.py} & $O(M \cdot n \cdot (n + m))$ no pior caso~\cite{kitsak2010} & \feito \\
\codigo{tabela\_de\_alcance} (Twitch PT) & \codigo{aplicacao/experimento.py} & o do anterior $+ O(n \log n)$ & \feito \\
\codigo{tau\_de\_kendall} (tau-b) & \codigo{aplicacao/influencia.py} & $O(n \log n)$~\cite{kendall1945,knight1966,christensen2005} & \feito \\
\codigo{comparar\_rankings} & \codigo{aplicacao/influencia.py} & $O(c \cdot r \cdot n \log n)$ & \feito \\
\codigo{tabela\_de\_concordancia} (Twitch PT) & \codigo{aplicacao/experimento.py} & $O(nm)$, pela intermediação & \feito \\
\codigo{beta\_recomendado} ($1{,}5\,\beta_c$) & \codigo{aplicacao/influencia.py} & $O(n)$ & \feito \\
\codigo{evolucao\_sir} (curvas S, I, R) & \codigo{aplicacao/influencia.py} & $O(M (n + m))$ no pior caso & \feito \\
\bottomrule
\end{tabularx}}

\subsection{Problema aplicado}
\pendente{a pergunta, a entrada, a saída e como validar --- texto da descrição formal da Etapa 1}

\subsection{Decisões e justificativas}
\begin{itemize}[leftmargin=*]
  \item SIR em tempo discreto, com $\beta$ um pouco acima do limiar epidêmico (0,0079 na Twitch PT),
        por exemplo 0,012, e $\gamma = 1$, seguindo~\cite{kitsak2010}. \pendente{confirmar os valores usados}
  \item O limiar é $\beta_c = \langle k \rangle / (\langle k^2 \rangle - \langle k \rangle)$~\cite{newman2002}:
        com $\gamma = 1$, $\beta$ é a chance de cada aresta transmitir, e o SIR vira percolação de
        arestas. Sem nó de grau 2 ou mais, o limiar é infinito. Só vale para grafo não direcionado.
  \item A simulação usa listas indexadas $0..n-1$ e um \codigo{bytearray} de estados que é zerado só
        nos nós atingidos: cada rodada custa o tamanho do surto, e não $O(n)$. Na Twitch PT, com
        $\beta = 0{,}012$ e $M = 100$, as 191.200 rodadas levam cerca de 22\,s.
  \item Um único gerador \codigo{random.Random(semente)} serve a todas as rodadas: com a mesma
        semente, o resultado se repete.
  \item O \codigo{simular\_sir} aceita $\gamma < 1$ (o infectado tenta de novo a cada passo), porque
        o caso de uso UC14 deixa o usuário escolher $\gamma$; o experimento usa $\gamma = 1$.
  \item O $\tau_b$ usa o algoritmo de Knight~\cite{knight1966}, na versão com empates descrita por
        \cite{christensen2005}: ordena os pares $(x, y)$, conta os empates e conta os pares
        discordantes como inversões de $y$ num merge sort. Custa $O(n \log n)$, e não $O(n^2)$:
        na Twitch PT, 4\,ms contra 0,7\,s da versão que compara os 1,8 milhão de pares. A versão
        direta fica nos testes, como referência.
  \item O $\tau_b$, e não o $\tau$ simples, porque o grau tem só 187 valores distintos e o núcleo só
        37: sem descontar os empates, as duas medidas sairiam prejudicadas.
  \item Se todos os valores de um lado são iguais, ou há menos de 2 nós, o $\tau_b$ é NaN: não há
        ordem para comparar.
\end{itemize}

\subsection{Corretude: testes e resultados}

{\small\begin{tabularx}{\linewidth}{@{}Y L{2.5cm} L{2.2cm} L{3.8cm}@{}}
\toprule
Teste & Base & Obtido & Referência \\
\midrule
SIR com $\beta = 0$ atinge só os iniciais; $\beta = \gamma = 1$ num conexo atinge todos & sintético, Karate & ok & --- \\
Alcance com $\beta = 1$ é o tamanho da componente (Hypothesis) & aleatórios & ok & --- \\
Alcance médio na aresta única, $(1+\beta)/2$, com $\beta = 0{,}3$ & sintético & 0,65 ($\pm$0,01) & 0,65 \\
Alcance médio na estrela ($n = 10$, $\beta = 0{,}4$), do centro e de uma folha & sintético & ok ($\pm$0,01) & $0{,}46$ e $0{,}268$ \\
Alcance com $\gamma = 0{,}5$, $\beta = 0{,}2$ na aresta única & sintético & ok ($\pm$0,01) & $2/3$ \\
Limiar: completo $1/(n-2)$, estrela $2/(n-2)$, ciclo 1 & sintético & ok & fórmulas \\
Kendall de rankings iguais e invertidos & sintético & 1 e $-1$ & 1 e $-1$ \\
Kendall com empate: $x = (1,2,2,3)$, $y = (1,2,3,4)$ & sintético & 0,9129 & $5/\sqrt{30}$ \\
Kendall no exemplo da documentação do SciPy (empates dos dois lados) & sintético & $-0{,}4714$ & $-0{,}4714$ \\
Kendall igual à definição direta, par a par (Hypothesis, 200 casos) & aleatórios & ok & --- \\
Limiar epidêmico & Twitch PT & 0,007876 & 0,0079 \\
$\tau$ entre grau e visualizações & Twitch PT & 0,4326 & $\approx$ 0,43 (NetworkX) \\
$\tau$ entre PageRank / k-core e visualizações & Twitch PT & 0,4386 / 0,4034 & $\approx$ 0,44 / 0,40 \\
\bottomrule
\end{tabularx}}

\subsection{Resultados}
Poder de espalhamento na Twitch PT (\codigo{python -m redes.aplicacao.experimento}; $\beta = 0{,}012$,
$\gamma = 1$, 100 rodadas por nó, semente 2026, cerca de 20\,s). Os atingidos por rodada vão de 1
(o surto morre na origem) a 88 (4,6\% da rede), com mediana 7,4. Os dez maiores alcances estão
todos no núcleo mais interno ($k_s = 37$); o nó 36 (grau 276) entra no top 10, e nós de grau 454 e
385 ficam de fora. A tabela completa está em \codigo{aplicacao/resultados/alcance\_twitch.csv}.

{\small\begin{tabularx}{\linewidth}{@{}r r r r r Y@{}}
\toprule
\# & Nó & Alcance & Grau & Núcleo & Visualizações \\
\midrule
1 & 127 & 4,60\% & 767 & 37 & 32.672.093 \\
2 & 467 & 4,53\% & 582 & 37 & 11.392.773 \\
3 & 1297 & 4,21\% & 587 & 37 & 7.724.802 \\
4 & 290 & 4,13\% & 590 & 37 & 116.392.808 \\
5 & 1476 & 4,01\% & 598 & 37 & 18.475.754 \\
6 & 1660 & 3,89\% & 475 & 37 & 33.101.623 \\
7 & 1758 & 3,88\% & 394 & 37 & 4.397.521 \\
8 & 1593 & 3,83\% & 368 & 37 & 3.094.236 \\
9 & 1320 & 3,82\% & 416 & 37 & 40.659.090 \\
10 & 36 & 3,69\% & 276 & 37 & 1.577.684 \\
\bottomrule
\end{tabularx}}

Concordância de cada centralidade com o alcance e com as visualizações ($\tau_b$ de Kendall; o
mesmo experimento, mais 30\,s para as centralidades). Arquivo:
\codigo{aplicacao/resultados/concordancia\_twitch.csv}.

{\small\begin{tabularx}{\linewidth}{@{}Y r r@{}}
\toprule
Centralidade & $\tau_b$ × alcance & $\tau_b$ × visualizações \\
\midrule
Grau & 0,7033 & 0,4326 \\
Proximidade & 0,7344 & 0,2889 \\
Harmônica & \textbf{0,7435} & 0,3031 \\
Intermediação & 0,5341 & 0,4133 \\
PageRank & 0,6592 & \textbf{0,4386} \\
Número de núcleo & 0,7301 & 0,4034 \\
\bottomrule
\end{tabularx}}

\begin{itemize}[leftmargin=*]
  \item \textbf{P1:} a harmônica é a que melhor prevê o alcance (0,7435), seguida da proximidade e
        do núcleo. Com as sementes 1, 2 e 3, a ordem das seis se repete. Com $\beta = 0{,}016$
        ($2\beta_c$), o núcleo passa à frente (0,795 contra 0,786 da harmônica).
  \item \textbf{P2:} o PageRank prevê melhor as visualizações (0,4386), quase empatado com o grau
        (0,4326). A proximidade e a harmônica, as melhores na P1, são as piores aqui.
  \item \textbf{H1} (núcleo supera o grau na P1): confirmada, 0,7301 contra 0,7033, e também com
        as outras sementes e com $\beta = 0{,}010$ e $0{,}016$.
  \item \textbf{H2} ($0 < \tau_b(c, w) < \tau_b(c, \rho)$ para toda $c$): confirmada. O próprio
        alcance concorda com as visualizações só com $\tau_b = 0{,}3376$: a audiência depende de
        mais coisas que a posição na rede de amizades.
\end{itemize}

\subsection{Limitações e pendências}
\begin{itemize}[leftmargin=*]
  \item \pendente{ex.: custo do SIR por Monte Carlo; número de rodadas usado}
\end{itemize}

\subsection{Registro}
\begin{description}[leftmargin=2.6cm,style=nextline]
  \registro{2026-11-19}{limiar epidêmico, SIR em tempo discreto e alcance médio por Monte Carlo
    em \codigo{aplicacao/influencia.py}, com 37 testes em \codigo{tests/test\_influencia.py}.}
  \registro{2026-11-19}{poder de espalhamento na Twitch PT: \codigo{aplicacao/experimento.py} grava
    o alcance de cada streamer em \codigo{aplicacao/resultados/alcance\_twitch.csv}; mais 6 testes.}
  \registro{2026-11-20}{$\tau_b$ de Kendall (Knight, $O(n \log n)$) e \codigo{comparar\_rankings};
    o experimento compara as 6 centralidades com o alcance e com as visualizações e grava
    \codigo{concordancia\_twitch.csv}; 58 testes no total.}
  \registro{2026-11-21}{integração com a interface: a tela de influenciadores usa
    \codigo{evolucao\_sir} (curvas e cores do mapa) e \codigo{poder\_de\_espalhamento} (alcance
    do ranking), com $\gamma = 1$, 100 rodadas e $\beta = 1{,}5\,\beta_c$ da rede ativa (0,012 na
    Twitch PT); a ``correlação de Pearson'' da tela, que tinha piso fixo (0,65) e valores fixos
    sem dados, virou o $\tau_b$ do grau, do núcleo e do PageRank contra o alcance e, na Twitch,
    contra as visualizações (com $\approx$ quando há amostra, em redes com mais de 2.000 nós);
    o pedido com $\beta$ ou $\gamma$ inválido é recusado sem mudar a sessão; 66 testes.}
\end{description}
